{
UK public procurement moves roughly \pounds300 billion a year, and when a losing
bidder challenges an award the resulting dispute is slow, expensive, and on the
evidence of the case law itself, usually settled long before a judge sees it.
This project asks whether interacting large language model agents can usefully
simulate such a dispute, and how their output should be evaluated when ground
truth is scarce and the model is capable of inventing evidence.

Three negotiating agents (a Contracting Authority, an Aggrieved Bidder, and a
Court modelled on the Technology and Construction Court), together with a
non-negotiating summariser, were built on LangGraph over a locally hosted Llama
3.1. Behaviour is specified through prompt engineering and structured context
injection rather than fine-tuning. The Court agent runs a conditional
verification protocol: it recomputes a scoring calculation exactly when the case
supplies both a formula and every input that formula needs, and otherwise
conducts a qualitative review in which writing down an unstated number is
prohibited. An extraction stage turns real judgments into structured scenarios,
and a rule-based pre-award risk screen addresses the industrial partner's actual
request, which was prevention rather than replay.

Evaluation covers 23 verified UK procurement judgments and 299 logged
negotiations. An audit for outcome leakage finds six merits scenarios that state
what the real court decided; on the eight that do not, the system matches the real
disposition in seven, against a majority-class floor of four. It needed structural
repair on 4 of 2,865 structured responses. Two findings matter more than that
headline. A single zero-shot prompt ties the full pipeline on directional accuracy
across the six cases where both were compared, so the architecture cannot be
defended on accuracy alone; and removing the Court agent deadlocks all 18 runs,
which is where it can be.
}
